\subsection{Completed simulation studies of PET design}
\label{sec:pet-campaigns}

The improvement and final-design studies tested whether better point estimates
could be combined with a calibrated uncertainty procedure. They improved
recovery substantially, but the terminal selection outcome is
\texttt{NO\_ELIGIBLE\_DESIGN}: no design or default was selected.
All results below concern signal-only simulation and method development,
not a real-data cross section or a publication uncertainty product
\cite{PETImprovement2026,PETFinalDesign2026}.

For unit-normalized histograms, recovery is
\begin{equation}
 R=1-\frac{\sum_b|u_b-t_b|}{\sum_b|p_b-t_b|},
 \label{eq:pet-campaign-recovery}
\end{equation}
where $u$, $t$ and $p$ are the unfolded, pseudodata-truth and prior
distributions on the stated reporting histogram. Thus $R=1$ is exact
recovery, $R=0$ retains the prior's discrepancy, and $R<0$ moves farther
from the target. A nearly unchanged marginal makes the denominator small;
joint hadron-content histograms and absolute residuals are needed in such
cases. This metric differs from the historical cellwise recovery proxy in
\S\ref{sec:gbdt-model-dependence}.

\paragraph{Improvement study: recovery gains did not qualify a design.}
Development, pilot, final and stress samples had separate roles. On twelve
fresh final draws, the historical PET control at three iterations recovered
$R=0.316$ of the seven-bin $\Eavail$ tilt. The frozen interventions at
$K^*=10$ gave the results in Table~\ref{tab:pet-improvement}.

\begin{table}[htbp]
\centering
\small
\begin{tabular}{@{}lrr@{}}
\toprule
Frozen intervention & Mean $R$ & One-sided 95\% lower bound \\
\midrule
A: iterations alone & 0.505 & 0.490 \\
B: reco summaries and one-hot truth PDG & 0.569 & 0.555 \\
C: efficiency-corrected truth step & 0.786 & 0.763 \\
\bottomrule
\end{tabular}
\caption{Improvement study at its declared ten-iteration operating point,
twelve final draws per configuration. B retains the carry-misses rule;
C trains the truth-step ratio on selected events and extrapolates it to
missed events. The unchanged aggregate adequacy floor is 0.556, with
additional regional mean requirements. B passes the point rule but its
lower bound is below the floor; A fails and C clears it. B versus A combines
two input changes and does not isolate the reco summaries.}
\label{tab:pet-improvement}
\end{table}

The 36 final and 36 stress runs completed. Coverage was not run: no frozen
candidate at $K^*=10$ met both adequacy and the stress requirement.
B moved away under neutron-content reweighting, and C under proton- and
neutron-content reweighting, despite these distortions being distinguishable
at reconstruction level. There were only two stress draws per case.
C at three iterations met the mean adequacy rule and had no moves-away
stress cell, but its lower bound was below the floor and that iteration was
not the frozen coverage point. It was not retrospectively promoted.
The combined detector-scale-plus-tilt case had no matching identifiability
test or unscaled control on the same draws; its recovery does not isolate
the response effect. The review left the PET failure mechanism unestablished.
Subsequent development showed that the predecessor's C neutron case barely
changed the scored marginal, whereas B made a genuine wrong-direction
push. These diagnostics motivated the next study without changing the
predecessor's frozen verdict.

\paragraph{Final design: representation, training and seed reproducibility.}
The next study used fixed development, final and sealed reserve banks.
Each unfolding had about 600,000 prior events and 600,000 pseudodata events;
final inference is conditional on those banks. Detector energy summaries
helped recover information omitted by the stored top-12-cluster cloud.
Categorical truth PDG encoding, a constant per-iteration learning rate and
longer truth-step training completed the hybrid designs. In development,
extending the truth step to 24 epochs reduced the fixed-event seed SD of
the primary recovery to 0.037 and 0.043, below the 0.05 limit. These are
seed-reproducibility measurements, not uncertainty coverage.

The frozen finalists were compact H2S1T24 at five iterations and larger
L128S1T24 at four, with approximately 47,000 and 970,000 detector-step
parameters. Pretrained PET2 improved on its scratch counterpart, but no
tested PET2 design passed all development screens: the stable annealed
variant failed the proton-topology screen. This is a result for the tested
recipes, not a general rejection of pretraining.

On 60 final-bank tilt draws, H2S1T24 and L128S1T24 reached
$R=0.895$ and $0.863$, respectively, against the historical PET control's
$0.306$. H2S1T24 passed the frozen accuracy, robustness and stability rules.
L128S1T24 failed B2, the neutron-down test on the nearly unchanged
$\Eavail$ marginal: mean residual minus
injected $L_1$ was $0.0122$, above $0.010$, versus H2S1T24's $0.0097$.
Both 95\% intervals cross the threshold, so the frozen point rule decides
eligibility but the statistical comparison establishes neither a robustness
ranking nor equivalence. The frozen B1 independence-based bound passes;
within-draw dependence prevents establishing its intended per-unit failure
probability bound. Nor does it show that the per-unit failure probability exceeds 0.10.
The compact design was not cheaper at the declared packing, so the
cost-based preference for it did not apply.

\paragraph{Coverage fails, with opposite behavior across acceptance regions.}
The H2S1T24 interval procedure uses six independent Poisson-bootstrap
members, each with its own training seeds. Its estimate is their mean
normalized histogram; intervals are
$\bar u\pm t_5 s\sqrt{1+1/6}$ at the stated confidence level, where $s$
is the member SD. On 120 final-bank development-tilt draws (720 unfoldings),
the results are those of Table~\ref{tab:pet-final-coverage}.

\begin{table}[htbp]
\centering
\small
\begin{tabular}{@{}lrr@{}}
\toprule
Scored region & Nominal 68\% coverage & Nominal 95\% coverage \\
\midrule
Aggregate $\Eavail$ & 0.896 & 0.990 \\
Low acceptance & 0.257 & 0.789 \\
Moderate acceptance & 0.701 & 0.964 \\
Good acceptance & 0.929 & 0.998 \\
\bottomrule
\end{tabular}
\caption{H2S1T24 K5 six-member interval coverage against each draw's own
pseudodata truth. Aggregate intervals over-cover: C1 fails decisively at
68\% (lower bound 0.858, ceiling 0.80). C4 also fails: mean 95\% half-widths
in six of seven bins are 1.32--1.72 times the allowed limit. Low acceptance
under-covers; no frozen rule gates that region (C3 covers only the moderate
and good regions). C2 and C3 pass.}
\label{tab:pet-final-coverage}
\end{table}

The same over- and under-coverage pattern survives comparison with the
final-bank population target. Its mechanism is not established; it is not
uniform conservatism. The proton-topology coverage stage C5 was skipped
after the decisive failure, and L128S1T24 coverage was never measured.
The terminal outcome remains \texttt{NO\_ELIGIBLE\_DESIGN}: H2S1T24 fails
interval calibration and L128S1T24 fails B2 by the frozen point rule (statistically
unresolved). Neither failure proves that the
point estimator is unusable. No repair or reserve-bank revalidation forms
part of these completed results.

\subsection{Exploratory PET--GBDT comparison on the same events}
\label{sec:pet-gbdt-paired}

The later comparison reused both finalists' point estimates and fitted a
scalar HistGradientBoosting OmniFold to the identical 352 final-bank units
per finalist, with the same events, weights, distortions, targets and scorer
\cite{PETGBDT2026}. Its efficiency-corrected miss treatment matches PET's;
detector inputs are muon kinematics, reconstructed $\Eavail$ and $q_3$,
stored-cluster energy sum and count. Truth inputs include four scalar
kinematics and species counts over the stored 12-token cloud, not full
particle multiplicities. The primary GBDT iteration, $k=7$, was chosen on
development data before final-bank use. This is a specific comparator,
distinct from the production scalar-5D LightGBM synthesis in
\S\ref{sec:gbdt-model-dependence}.

PET recovers both $\Eavail$ tilts better than this GBDT
(Fig.~\ref{fig:pet-gbdt-paired}). On the development tilt, paired gains in
$R$ are $+0.102$ [0.084, 0.119] for H2S1T24 and $+0.070$ [0.051, 0.089]
for L128S1T24 over 60 draws. Summed per-bin mean squared residuals are
$(3.33,5.26,9.78)\times10^{-4}$ for H2S1T24, L128S1T24 and GBDT.
Residuals compare each estimate to its own draw's target; their spread is
not unconditional estimator variance. On the proton-topology endpoint,
H2S1T24's gain at GBDT $k=7$ is $+0.055$ [0.035, 0.075], but at
$k=10$ it is $+0.009$ [$-0.012$, 0.030]; L128S1T24 is level at $k=7$
($+0.012$ [$-0.007$, 0.031]) and below the GBDT at $k=10$ ($-0.034$
[$-0.054$, $-0.015$]). The $\Eavail\times q_3$
endpoint does not resolve a difference.

\begin{figure}[htbp]
\centering
\includegraphics[width=\linewidth]{pet_gbdt_paired_differences.pdf}
\caption{Existing exact-paired comparison of the PET finalists with GBDT
at its development-chosen $k=7$. Positive differences favor PET.
E0 and E3 are the two $\Eavail$ tilts (60 draws); E1 subdivides E0 by
acceptance; E4 is $\Eavail\times$ proton class and E5 is
$\Eavail\times q_3$ (40 draws). Bars are unadjusted two-sided 95\%
Student-$t$ intervals, conditional on the frozen banks. Exploratory point
estimates only: no uncertainty comparison or finalist ranking follows.}
\label{fig:pet-gbdt-paired}
\end{figure}

The broader distortion library reverses some of this ordering
(Fig.~\ref{fig:pet-gbdt-library}). GBDT is better on all three tested
generator-model reweightings (NuWro, NuWro$'$ and GiBUU) by 0.12--0.25 in
$R$, and on the narrow energy bump. These are fixed-response simulation
reweightings, not independent detector simulations or real-data comparisons.
PET leads on several species distortions. Its raw muon-scale result largely
includes its underlying tilt gain: the response-specific paired contrast
is only $+0.016$ [$-0.005$, 0.037] for H2S1T24 and $+0.028$
[0.013, 0.043] for L128S1T24. Signed response contrasts alone are not
robustness rankings.

\begin{figure}[H]
\centering
\includegraphics[width=0.94\linewidth]{pet_gbdt_library_paired.pdf}
\caption{Paired recovery differences across the existing distortion library,
GBDT at $k=7$. Each case uses its natural histogram; most have eight
draws, with D4c proton-up and D3 $+0.35$ ($\Eavail\times q_3$) using 40.
Intervals are unadjusted 95\% Student-$t$ intervals conditional on the
banks. The response-plus-tilt rows include the tilt contribution and do
not isolate detector-response sensitivity. The generator rows show that
PET's tilt advantage does not extend to every tested departure.}
\label{fig:pet-gbdt-library}
\end{figure}

All contrasts are exploratory: the PET final-bank results had already been
inspected, intervals are not adjusted for multiple comparisons, and inference
conditions on the finite banks at study scale (about 600,000 pseudodata
events per unfolding). On the two matched development draws (descriptive),
binned iterative Bayesian unfolding and AUSSIE with kinematic truth inputs
match or beat PET on the one-dimensional tilt but fail on hadron-content or
generator cases; the observed lead is over this HGB OmniFold, not scalar
unfolding in general. There is no GBDT interval procedure in this comparison,
and the measured PET procedure failed calibration. Thus no uncertainty
comparison, universal method ordering, finalist ranking or adoption follows.
The completed comparison leaves \texttt{NO\_ELIGIBLE\_DESIGN} unchanged.

The improvement and final-design numerical conclusions have the committed
independent checks recorded with their reports. The paired comparison's
numbers also reproduce independently; its post-second-review wording and
prospective design repairs were not re-reviewed. These distinctions do not
upgrade any result to publication uncertainty status. Ledger entries
VL164--VL168 record the final-design and comparison results; the cited
reports route the frozen decisions and closed review dispositions.
